\section{Experimental Setup}\label{sec:experimental_setup}

\subsection{CNC Machine}

All experiments are conducted on a Bantam Tools Explorer Desktop CNC Milling Machine, a compact 3-axis vertical milling center with a \SI{26000}{RPM} brushless DC spindle and a work envelope of $140 \times 114 \times \SI{38}{\milli\meter}$. The machine uses stepper-driven ball screws on the X, Y, and Z axes. This desktop platform differs from industrial CNC mills, which use servo drives with closed-loop position feedback, have orders-of-magnitude higher structural stiffness, and operate at higher cutting forces. It nevertheless provides a controlled, accessible testbed for developing and validating the air-cut referencing methodology. Generalization to industrial platforms remains an open question addressed in Section~\ref{sec:discussion}.

A single tool is used throughout: a 2-flute high-speed steel flat end mill, $\sfrac{1}{4}$\,inch (\SI{6.35}{\milli\meter}) diameter (T2). Spindle speed is fixed at $S = \SI{12000}{RPM}$ for all runs. At this speed, the tooth-passing frequency is \SI{400}{\hertz} ($2 \times 200$~rev/s), well above the sensor system's Nyquist limit (Section~\ref{sec:sampling}).

The workpiece material is TIVAR ultra-high molecular weight polyethylene (UHMW-PE), chosen for consistent cutting behavior, minimal tool wear, and safe data collection without coolant or hazardous chips.

\subsection{Toolpath Types}

Three toolpath types generated by Autodesk Fusion~360 are studied:

\begin{itemize}
    \item \textbf{Adaptive clearing:} A high-efficiency roughing strategy that dynamically adjusts the stepover to maintain approximately constant tool engagement, following a trochoidal-like spiral path.
    \item \textbf{Facing:} A surface-finishing operation with back-and-forth linear passes at a fixed stepover.
    \item \textbf{Pocketing:} A material removal strategy for enclosed regions using contour-parallel passes to clear a rectangular pocket.
\end{itemize}

Active cutting uses a \SI{0.150}{in} (\SI{3.81}{\milli\meter}) stepover and a \SI{0.025}{in} (\SI{0.635}{\milli\meter}) depth of cut.

\subsection{Sensor Instrumentation}

The sensing system comprises six Arduino-based IMU modules and a 4-channel electrical current measurement system, at a total hardware cost of approximately \$492~USD.

\subsubsection{IMU Modules}

Each IMU module is an \textbf{Arduino Nano 33 BLE Sense Lite} board (\$33~USD each), which integrates the following sensors:

\begin{itemize}
    \item \textbf{LSM9DS1} 9-axis IMU (STMicroelectronics): 3-axis accelerometer ($\pm 2$g to $\pm 16$g, 16-bit resolution, $\sim$0.061~mg/LSB at $\pm 2$g range), 3-axis gyroscope ($\pm 245$ to $\pm 2000$~dps), 3-axis magnetometer ($\pm 400$ to $\pm 1600$~$\mu$T). Native sampling rate: \SI{119}{\hertz} (Nyquist: \SI{59.5}{\hertz}).
    \item \textbf{LPS22HB} barometric pressure and temperature sensor (STMicroelectronics): Measures atmospheric pressure (\SI{260}{\hecto\pascal} to \SI{1260}{\hecto\pascal}) and die temperature. The pressure channel measures barometric conditions, not cutting-related pressure; it varies by collection day ($\sim$100.0 vs.\ $\sim$101.5~kPa) and acts as a session-level confound rather than a process indicator.
    \item \textbf{APDS-9960} gesture, proximity, and RGB color sensor (Broadcom): The proximity channel ($0$--$255$) detects nearby objects via IR reflection; the RGBA color channels measure ambient light and reflected LED intensity. In this deployment, proximity readings are near-constant (determined by sensor housing geometry), and color channels encode LED shadow patterns that weakly correlate with gantry position.
    \item \textbf{MP34DT05} PDM microphone (STMicroelectronics): Acoustic RMS amplitude is computed onboard before transmission, yielding a single-value vibration/acoustic energy proxy per sample.
\end{itemize}

This produces 17 channels per module: $a_x, a_y, a_z$ (accelerometer), $\omega_x, \omega_y, \omega_z$ (gyroscope), $m_x, m_y, m_z$ (magnetometer), pressure, temperature, proximity, $R, G, B, A$ (color), and RMS (acoustic). With six modules, the IMU subsystem provides $6 \times 17 = 102$ channels.

\textbf{Mounting:} The Arduino boards are soldered onto breadboards, which are affixed to machine surfaces using cyanoacrylate adhesive for rigid structural coupling. This provides adequate mechanical transmission for the low-frequency regime ($< 60$~Hz) captured by the sensor system, though it would be insufficient for high-frequency vibration monitoring on industrial machines.

\textbf{Sensor placement:} Six modules are deployed across three structural groups (Table~\ref{tab:sensor_placement}), selected from 16 initially deployed boards based on $\geq 93\%$ data availability (10 boards excluded due to intermittent BLE connectivity dropouts).

\begin{table}[H]
\caption{Sensor module locations and structural groups.}\label{tab:sensor_placement}
\centering
\begin{tabular}{llll}
\toprule
\textbf{Module ID} & \textbf{Location} & \textbf{Structural Group} & \textbf{Rationale} \\
\midrule
frame\_l2  & Frame, left mid-height  & Frame (3 sensors) & Captures frame vibration \\
frame\_l3  & Frame, left upper       & Frame              & Vertical vibration path \\
frame\_r2  & Frame, right mid-height & Frame              & Bilateral symmetry \\
spindle2   & Spindle housing         & Spindle (1 sensor) & Closest to cutting zone \\
y\_bed\_\_3 & Moving table, pos.\ 3  & Bed (2 sensors)    & Near workpiece \\
y\_bed\_\_4 & Moving table, pos.\ 4  & Bed                & Bed vibration redundancy \\
\bottomrule
\end{tabular}
\end{table}

\subsubsection{Electrical Current Sensors}

Four \textbf{BAYITE Hall-effect current sensors} connected to a \textbf{Measurement Computing (MCC) USB-1608G} data acquisition module measure motor current draw:

\begin{itemize}
    \item Channel 0: Spindle motor current (proportional to cutting torque at steady-state RPM)
    \item Channels 1--3: X, Y, Z axis stepper motor current
\end{itemize}

The DAQ samples at \SI{1}{\kilo\hertz} in single-ended mode ($\pm 10$~V range). Each motor channel produces two derived features: instantaneous current and current acceleration (numerical derivative), yielding $4 \times 2 = 8$ electrical channels total.

\subsubsection{Data Synchronization and Effective Sampling Rate}\label{sec:sampling}

All sensor streams are aggregated on a \textbf{Raspberry Pi 4}. IMU modules communicate via Bluetooth Low Energy (BLE) wireless links; the MCC DAQ connects via USB. The streams are resampled and synchronized to the G-code execution timeline, yielding an effective aligned sampling rate of approximately \SI{4}{\hertz} ($\Delta t \approx \SI{0.25}{\second}$). This common rate is set by the slowest reliable cross-stream synchronization achievable over the BLE wireless links and the G-code-timeline alignment, not by the native sensor bandwidths (\SI{119}{\hertz} IMU, \SI{1}{\kilo\hertz} DAQ); retaining the native high-rate streams for frequency-domain analysis is left to future work. Each timestep produces a feature vector of $C = 102 + 8 = 110$ channels.

\textbf{Bandwidth limitations.} The \SI{4}{\hertz} effective rate severely limits the frequency content available for classification. The tooth-passing frequency (\SI{400}{\hertz}), chatter frequencies (typically \SI{500}{}--\SI{5000}{\hertz}), and stepper motor electrical frequencies (\SI{50}{}--\SI{500}{\hertz}) are all far above the \SI{2}{\hertz} Nyquist limit. Consequently, the sensor data captures \emph{quasi-static machine motion} rather than cutting-induced vibration dynamics: trajectory envelopes, low-frequency structural modes ($< 2$~Hz), and slowly-varying thermal/magnetic field changes. The temporal binning approach (Section~\ref{sec:methodology}) is designed to extract discriminative information from these quasi-static signals.

\textbf{Decimation and aliasing.} The resampling from native rates (\SI{119}{\hertz} for IMU, \SI{1}{\kilo\hertz} for DAQ) to the \SI{4}{\hertz} aligned rate discards all frequency content above \SI{2}{\hertz}. No explicit anti-aliasing filter is applied before decimation; however, the BLE wireless transmission introduces variable latency that acts as an uncontrolled low-pass filter, and the on-board firmware averaging (particularly for the MP34DT05 RMS computation) provides partial anti-aliasing for the acoustic channel. The \SI{1}{\kilo\hertz} DAQ captures motor current waveforms at adequate bandwidth (stepper electrical frequencies $\leq$\SI{500}{\hertz}), but this information is lost after decimation to \SI{4}{\hertz}. Future work should retain the high-rate DAQ data for frequency-domain analysis.

\textbf{Sensor calibration.} The LSM9DS1 IMU uses factory-trimmed calibration. No additional offset or gain calibration was performed for this deployment. Typical MEMS accelerometer bias stability is \SI{10}{}--\SI{50}{\milli\g}, which is comparable to the inter-type acceleration differences observed ($\sim$\SI{20}{\milli\g}; see Table~\ref{tab:temporal_cv}). This means the smallest observed differences may be near the sensor's bias uncertainty floor, a limitation that higher-grade MEMS sensors would address.

\subsection{Dataset Summary}

Table~\ref{tab:dataset} summarizes the dataset composition.

\begin{table}[H]
\caption{Dataset composition: number of runs per condition and toolpath type.}\label{tab:dataset}
\centering
\begin{tabular}{lccc}
\toprule
\textbf{Toolpath Type} & \textbf{Air-Cut Runs} & \textbf{Active-Cutting Runs} & \textbf{Total} \\
\midrule
Adaptive clearing & 20 & 17 & 37 \\
Facing            & 16 & 15 & 31 \\
Pocketing         & 19 & 19 & 38 \\
\midrule
\textbf{Total}    & 55 & 51 & 106 \\
\bottomrule
\end{tabular}
\end{table}

The dataset comprises 106 CSV files: 55 air-cut and 51 active-cutting runs. An additional 14 damaged-spindle air-cut runs (5 adaptive / 4 facing / 5 pocketing, collected in a single session with one spindle drive band removed, same instrumentation) are used only in the damage-detection experiment (Section~\ref{sec:damage}) and are excluded from all other analyses. Run durations range from $\sim$\SI{75}{\second} (pocketing) to $\sim$\SI{190}{\second} (facing), yielding 300--775 timesteps per run at \SI{4}{\hertz}. The air-cut and active-cutting conditions used different CAM-generated G-code programs (different stepovers, feed rates, and coordinates), preventing direct G-code-aligned signal subtraction and motivating the temporal binning approach (Section~\ref{sec:methodology}).
